\section{Unified Mobile Shopping, Booking and Service Actions}
\label{sec:pilot-unified-mobile-service-actions}

The unified Pilot mobile application now exposes shopping, booking, route and music actions through
one governed private service surface.  This is a proposal and approval system, not an autonomous
checkout agent.  Research or shared-television conversation may begin an action, but a named
person's trusted phone must complete any missing private details, inspect the exact scope and
approve it before an authorized provider adapter can run.

\subsection{User Experience and Supported Actions}

The Personal space includes a Shopping shortcut that opens the combined service-action centre.
The person can prepare a product search or basket request, a reservation, a destination route or a
music request.  The interface shows the service, intended action, exact parameters, missing fields
and current state.  A prepared request has no external effect.

The lifecycle is explicit:
\begin{enumerate}
  \item \texttt{needs\_details}: required information is missing and must be supplied privately;
  \item \texttt{awaiting\_private\_approval}: the complete immutable scope is ready for review;
  \item \texttt{approved\_pending\_adapter}: the person approved it, but it has not executed;
  \item \texttt{executing}: one authorized provider attempt may be in flight;
  \item \texttt{execution\_failed\_unknown\_effect}: an uncertain provider result requires reconciliation;
  \item \texttt{executed}: an adapter returned a verified result and Pilot recorded a receipt; or
  \item \texttt{rejected} or \texttt{cancelled}: no provider action is permitted.
\end{enumerate}

Approval and execution are separate controls.  If the selected person has no authorized account
binding or no adapter is active, Pilot preserves the approved proposal and reports the missing
connection.  It does not open an uncontrolled checkout page or claim success.

\subsection{Private-Identity and Legacy-Service Bridge}

The earlier Service Execution Hub used a household service persona, while the production mobile
application uses possession-bound private identities.  A one-to-one bridge now resolves these
without combining household members.  When a mother has exactly one active adult identity, that
adult may adopt the legacy local-owner service record, preserving earlier prepared actions and
opaque service bindings.  A multi-adult household never assigns that legacy account by arrival
order: each adult and child receives an isolated service persona, routing key and provider bindings
until the owner explicitly connects an account.  The internal service-persona identifier is
replaced by the production private identity in every phone projection.

This bridge preserves the principle that a shared television cannot decide whose card, booking
account, map history or media account should be used.  The trusted phone determines the person;
the mother brain selects only that person's authorized binding.  Child profiles may prepare a
request for discussion, but cannot approve or execute an external service effect without a future
guardian workflow.

\subsection{Signed Capability Boundary}

The phone uses four narrow capability scopes:
\begin{itemize}
  \item \texttt{services.read} for the person's proposals, connection state and receipts;
  \item \texttt{services.propose} for preparation and missing-detail completion;
  \item \texttt{services.approve} for accepting or rejecting one unchanged parameter hash; and
  \item \texttt{services.execute} for one previously approved provider operation.
\end{itemize}

Every request is signed by the phone's non-exportable Ed25519 key and validated against its mother-
signed certificate, active persona and short-lived capability lease.  A persona substitution or
changed parameter hash fails closed.  Preparation uses an idempotency key bound to the canonical
service, action and parameter hash; an identical retry returns the original proposal, while reuse
for altered details is rejected.

\subsection{Credential and Payment Safety}

Provider credentials remain behind opaque \texttt{vault://} or \texttt{provider://} references on
the mother brain.  They are never returned to the phone, television or device mesh.  Recursive
parameter validation rejects raw passwords, secrets, tokens, card numbers, CVVs, CVCs, PINs and
private keys.  The interface may record a budget or product requirement, but never payment-card
data.  Payment selection and regulated transaction handling require a separately authorized
provider flow and are not implied by this implementation.

\subsection{Execution and Verified Receipts}

Execution is permitted only when the stored proposal belongs to the same mapped service persona,
the same private identity approved it and the supplied hash still matches the stored parameters.
The Service Execution Hub then resolves that person's opaque provider binding and calls the
registered adapter.  The adapter must return a verified result.  Pilot stores a minimal normalized
receipt containing the proposal, service, action, external reference, result hash and timestamp;
the raw provider response is not retained.

If an adapter fails after an attempt starts, Pilot records an unknown-effect state and forbids an
automatic retry until reconciliation.  A completed proposal returns its existing receipt on repeat
execution, preventing duplicate purchasing or booking caused by a network retry.

\subsection{Verification and Remaining Qualification}

Automated tests cover production-to-household identity bridging, preserved service bindings,
persona isolation, idempotent preparation, immutable-hash approval, raw payment-secret rejection,
verified adapter receipts and private projection.  Surface tests cover the Shopping shortcut,
dynamic action forms, exact proposal review, missing-detail completion, separate approval and
execution controls, and verified receipt presentation.  At coded closure of \texttt{MOB-0507}, all
463 AION Fabric tests passed and the optimized Next.js build statically generated the mobile route.

No real purchase, booking, map-account change or music-account action was made during verification.
Production qualification still requires explicit user authorization and a verified receipt from
each real provider.  Those external qualifications remain open and cannot be closed by a mock
adapter or interface demonstration.
